% ============================================================
% Page 2 — Alban aux temps « gallo-romains »
% Transcription fidèle du document source.
% Mise en page volontairement simple : typographie à définir
% ultérieurement dans le préambule général.
% ============================================================

\thispagestyle{plain}

\vspace*{12mm}

\begin{center}
    {\Large\bfseries Alban aux temps « gallo-romains »}

    \vspace{2mm}
    {\itshape (800 ans avant J.C. -- 400 ans après)}
\end{center}

\vspace{18mm}

\noindent
Des siècles durant, des populations celtiques et romaines se heurtent,
se croisent, cohabitent, finalement s'entendent. Les Gaulois arrivent
jusqu'en Asie mineure (Turquie aujourd'hui) : qu'on se souvienne de la
lettre de Saint Paul aux Galates. Les Romains entraînent leurs légions
jusqu'en Angleterre\ldots{} Chez nous, des débris de poterie, des lampes,
des monnaies attestent la présence de Rome. La Société des Arts et Belles
Lettres du Tarn recevait le 28 juin 1895 de Q. Soulage des lampes romaines
provenant des mines d'Alban. Le 12 mai 1905, on lui présentait une pièce
d'or à l'effigie de Sévère. Le 19 mars 1952, au domaine de Peyre-Augé,
était exhumée\textsuperscript{1} une amphore de forme sphérique du
2ème siècle. Une pièce semblable est au musée de Tours, une autre à
Mende. Le vase gallo-romain « garum » a été plaisamment baptisé
« marmite d'Astérix » : on peut le voir au fond de l'église d'Alban.
Mais à partir de quand Alban « sent-il » vraiment le romain ?

\vfill

\noindent\rule{55mm}{0.4pt}

\vspace{1mm}
{\small\itshape Fouilles de M. Roustan}
